\subsection{MATRICES OVER A COMMUTATIVE GROUP}

Let G be a commutative group (written additively). The set of matrices over G, with the given indexing sets I, K, has a commutative group structure since it is the set of mappings from \(I \times K\) to G; this group is written additively, so that if \(M = (m_{\iota\kappa})\) and \(M' = (m_{\iota\kappa}')\) are two of its elements, then

\[
M + M ^ {\prime} = (m _ {\iota \kappa} + m _ {\iota \kappa} ^ {\prime});
\]

the identity element of this group is therefore the matrix all of whose elements are zero (called the zero matrix). Clearly

\[
{ } ^ { t } ( M + M ^ { \prime } ) = { } ^ { t } M + { } ^ { t } M ^ { \prime } .\tag{2}
\]

The sum of two matrices is thus only defined if the indexing sets of the rows and the columns are the same for the two matrices.

Let \(\mathbf{H}'\) , \(\mathbf{H}''\) be two sets, \(\mathbf{G}\) a commutative group (written additively) and \(f: (h', h'') \mapsto h'h''\) a mapping from \(\mathbf{H}' \times \mathbf{H}''\) to \(\mathbf{G}\) . Given two matrices

\[
M ^ {\prime} = \left(m _ {i k} ^ {\prime}\right) _ {(i, k) \in I \times K}, \quad M ^ {\prime \prime} = \left(m _ {k l} ^ {\prime \prime}\right) _ {(k, l) \in K \times L}
\]

over H' and H'' respectively such that the indexing set K of the columns of M' is finite and equal to the indexing set of the rows of M'', the product of M' and M'' via f, denoted by M'M'' or f(M', M''), is the matrix

\[
\left(\sum_ {k \in \mathbf {K}} m _ {i k} ^ {\prime} m _ {k l} ^ {\prime \prime}\right) _ {(i, l) \in \mathbf {I} \times \mathbf {L}}\tag{3}
\]

over G.

The above definition supposes that the indexing set of the columns of \(M'\) is equal to the indexing set of the rows of \(M''\) ; in particular the product \(M''M'\) has no meaning if \(I \neq L\) . In formula (3) the elements of the same

row of \(M'\) figure multiplied on the right by the elements of the same column of \(M''\) ; the multiplication is said to be made ``rows by columns''.

Let \(f^0\) be the mapping \((h'', h') \mapsto h'h''\) of \(H'' \times H'\) to G; it follows immediately from the definitions that

\[
{ } ^ { t } ( M ^ { \prime } M ^ { \prime \prime } ) = { } ^ { t } M ^ { \prime \prime } . { } ^ { t } M ^ { \prime }\tag{4}
\]

where the product on the left (resp. right) hand side is calculated via \(f\) (resp. via \(f^0\) ).

When \(\mathbf{H}'\) and \(\mathbf{H}''\) are themselves commutative groups (written additively) and \(f\) is \(\mathbf{Z}\) -bilinear ( \(\S 3\) , no. 1), the distributivity formulae

\[
\left\{ \begin{array}{l} (M ^ {\prime} + N ^ {\prime}) M ^ {\prime \prime} = M ^ {\prime} M ^ {\prime \prime} + N ^ {\prime} M ^ {\prime \prime} \\ M ^ {\prime} (M ^ {\prime \prime} + N ^ {\prime \prime}) = M ^ {\prime} M ^ {\prime \prime} + M ^ {\prime} N ^ {\prime \prime} \end{array} \right.\tag{5}
\]

are immediately verified, the indexing sets being such that the sums and products appearing are defined.

Now let \(H_{1}\) , \(H_{2}\) , \(H_{3}\) , \(H_{12}\) , \(H_{23}\) and H be commutative groups (written additively), \(f_{12}:H_{1} \times H_{2} \to H_{12}\) , \(f_{23}:H_{2} \times H_{3} \to H_{23}\) mappings and

\[
f _ {3}: \mathrm{H} _ {1 2} \times \mathrm{H} _ {3} \rightarrow \mathrm{H}, \quad f _ {1}: \mathrm{H} _ {1} \times \mathrm{H} _ {2 3} \rightarrow \mathrm{H}
\]

Z-bilinear mappings; suppose further that, for all \(x_{i} \in H_{i}\) (i = 1, 2, 3)

\[
f _ {3} (f _ {1 2} (x _ {1}, x _ {2}), x _ {3}) = f _ {1} (x _ {1}, f _ {2 3} (x _ {2}, x _ {3}))
\]

(which may also be written as above \((x_{1}x_{2})x_{3} = x_{1}(x_{2}x_{3}))\) ; then, if \(M' = (m'_{rs})\) , \(M'' = (m''_{st})\) , \(M''' = (m''_{tu})\) are matrices over \(\mathrm{H}_1\) , \(\mathrm{H}_2\) , \(\mathrm{H}_3\) respectively,

\[
(M ^ {\prime} M ^ {\prime \prime}) M ^ {\prime \prime \prime} = M ^ {\prime} (M ^ {\prime \prime} M ^ {\prime \prime \prime})\tag{6}
\]

when the products on the two sides (calculated respectively via \(f_{12}\) , \(f_{3}\) , \(f_{23}\) and \(f_{1}\) ) are defined; for

\[
\begin{array}{r l} \sum_ {t} \left(\sum_ {s} m _ {r s} ^ {\prime} m _ {s t} ^ {\prime \prime}\right) m _ {t u} ^ {\prime \prime \prime} & = \sum_ {t} \sum_ {s} (m _ {r s} ^ {\prime} m _ {s t} ^ {\prime \prime}) m _ {t u} ^ {\prime \prime \prime} = \sum_ {s} \sum_ {t} m _ {r s} ^ {\prime} (m _ {s t} ^ {\prime \prime} m _ {t u} ^ {\prime \prime \prime}) \\ & = \sum_ {s} m _ {r s} ^ {\prime} \left(\sum_ {t} m _ {s t} ^ {\prime \prime} m _ {t u} ^ {\prime \prime \prime}\right) \end{array}
\]

by virtue of the hypotheses made.

The two sides of (6) are also denoted by \(M'M''M'''\) . Analogous conventions are made for products of more than three factors.

Remark. The above formulae extend to a more general situation. To be precise:

\begin{enumerate}
\def\labelenumi{(\alph{enumi})}
\item
  Suppose \(H = \bigcup_{(\iota, \kappa) \in I \times K} G_{\iota\kappa}\) where each \(G_{\iota\kappa}\) is a commutative group written additively; then the sum \(M + M'\) may be defined when, for each ordered pair \((\iota, \kappa)\) , \(m_{\iota\kappa} \in G_{\iota\kappa}\) and \(m'_{\iota\kappa} \in G_{\iota\kappa}\) .
\item
  Let I, K, L be three sets with K assumed finite and let \(H' = \bigcup_{(i, k) \in I \times K} H'_{ik}\) , \(H'' = \bigcup_{(k, l) \in K \times L} H''_{kl}\) , \(H = \bigcup_{(i, l) \in I \times L} H_{il}\) be three sets; suppose that each \(H_{il}\) is a commutative group written additively and for each triple \((i, k, l)\) let
\end{enumerate}

\[
f _ {i k l}: \mathrm{H} _ {i k} ^ {\prime} \times \mathrm{H} _ {k l} ^ {\prime \prime} \rightarrow \mathrm{H} _ {i l}
\]

be a mapping. Then if \(M' = (m_{ik}')_{(i,k) \in I \times K}\) , \(M'' = (m_{kl}'')_{(k,l) \in K \times L}\) are matrices such that \(m_{ik}' \in H_{ik}'\) and \(m_{kl}'' \in H_{kl}''\) for all \(i, k, l\) we can define the product \(M'M''\) via the \(f_{ikl}\) . We leave to the reader the task of writing down and proving the formulae analogous to (4), (5) and (6).

